% GENERATED FILE. Edit canonical lessons, then run npm run generate:books.
% canonical-source-hash: fnv1a64:59503941d61a3d29
% canonical-lessons: JA-C333-wakamono, JA-C333-koureisha, JA-C333-jinkou, JA-C333-kaigo-suru, JA-C333-kosodate

\chapter{Young and Old}
\label{ch:ja-young-and-old}
\hlchaptermodality{333}

\begin{chapteropening}
\textbf{By the end of this chapter:} \emph{I can say how the young and the old live, and why.}

Complete the last lesson of chapter 333: I can say how the young and the old live, and why.
\end{chapteropening}

\section[\texorpdfstring{\ja{わかもの} (wakamono)}{わかもの (wakamono)}]{\ja{わかもの} (wakamono) — young people; a young person}

\label{lesson:JA-C333-wakamono}



\begin{quote}
Before the new one: say the Japanese for discrimination, then the Japanese for all members.
\end{quote}

\subsection*{You'll want to know: \ja{わかもの}}
\textbf{\ja{わかもの}} — \emph{wakamono} — \textquotedblleft{}young people; a young person\textquotedblright{}.

\begin{quote}
\textbf{\ja{いなかには} \ja{わかものが} \ja{すくないと} \ja{おもいます。}} — \emph{inaka ni wa wakamono ga sukunai to omoimasu} — \textquotedblleft{}I think there are few young people in the countryside.\textquotedblright{}
\end{quote}

\begin{cousinweb}
\textbf{\ja{わかもの}} (\emph{wakamono}) is a native Japanese word (wago), inherited rather than borrowed from Chinese: it is built from \emph{waka}, \textquotedblleft{}young (wakai)\textquotedblright{} and \emph{mono}, \textquotedblleft{}person\textquotedblright{}, and its kanji were attached later as kun'yomi (JMdict).
\end{cousinweb}

\begin{grammarlens}[title={What you've built}]
One more word for giving reasons and opinions.
\end{grammarlens}

\subsection*{Your turn}
\begin{itemize}
\raggedright
  \item \emph{Say it:} \emph{wakamono}
  \item \emph{Say it:} \emph{wakamono}, once more
  \item \emph{Say it:} \emph{inaka ni wa wakamono ga sukunai to omoimasu} — I think there are few young people in the countryside.
  \item \emph{Recall:} say the Japanese for discrimination, then the Japanese for all members, then say \emph{wakamono} again
\end{itemize}

\begin{tcolorbox}[breakable,colback=teal!4,colframe=teal!35!black,title={Before you move on}]
What does \ja{わかもの} mean? (\textquotedblleft{}Young people; a young person\textquotedblright{}.) Say it once more.
\end{tcolorbox}
\section[\texorpdfstring{\ja{こうれいしゃ} (kōreisha)}{こうれいしゃ (kōreisha)}]{\ja{こうれいしゃ} (kōreisha) — an elderly person}

\label{lesson:JA-C333-koureisha}



\begin{quote}
Before the new one: say the Japanese for all members, then the Japanese for young people.
\end{quote}

\subsection*{You'll want to know: \ja{こうれいしゃ}}
\textbf{\ja{こうれいしゃ}} — \emph{kōreisha} — \textquotedblleft{}an elderly person\textquotedblright{}.

\begin{quote}
\textbf{\ja{こうれいしゃが} \ja{ふえたと} \ja{おもいます。}} — \emph{kōreisha ga fueta to omoimasu} — \textquotedblleft{}I think the number of elderly people has grown.\textquotedblright{}
\end{quote}

\begin{cousinweb}
\textbf{\ja{こうれいしゃ}} (\emph{kōreisha}) is Sino-Japanese (kango): three Chinese characters read with their on'yomi, readings borrowed from Chinese. They are \emph{kō}, \textquotedblleft{}high\textquotedblright{} (Middle Chinese \emph{kaw}), \emph{rei}, \textquotedblleft{}age\textquotedblright{} (\emph{leng}) and \emph{sha}, \textquotedblleft{}someone\textquotedblright{} (\emph{tsyaeX}) (KANJIDIC2; Middle Chinese in Baxter's transcription).
\end{cousinweb}

\begin{grammarlens}[title={What you've built}]
One more word for giving reasons and opinions.
\end{grammarlens}

\subsection*{Your turn}
\begin{itemize}
\raggedright
  \item \emph{Say it:} \emph{kōreisha}
  \item \emph{Say it:} \emph{kōreisha}, once more
  \item \emph{Say it:} \emph{kōreisha ga fueta to omoimasu} — I think the number of elderly people has grown.
  \item \emph{Recall:} say the Japanese for all members, then the Japanese for young people, then say \emph{kōreisha} again
\end{itemize}

\begin{tcolorbox}[breakable,colback=teal!4,colframe=teal!35!black,title={Before you move on}]
What does \ja{こうれいしゃ} mean? (\textquotedblleft{}An elderly person\textquotedblright{}.) Say it once more.
\end{tcolorbox}
\section[\texorpdfstring{\ja{じんこう} (jinkō)}{じんこう (jinkō)}]{\ja{じんこう} (jinkō) — population}

\label{lesson:JA-C333-jinkou}



\begin{quote}
Before the new one: say the Japanese for young people, then the Japanese for an elderly person.
\end{quote}

\subsection*{You'll want to know: \ja{じんこう}}
\textbf{\ja{じんこう}} — \emph{jinkō} — \textquotedblleft{}population\textquotedblright{}.

\begin{quote}
\textbf{\ja{まちの} \ja{じんこうが} \ja{すくなく} \ja{なったので、みせが} \ja{しまりました。}} — \emph{machi no jinkō ga sukunaku natta node, mise ga shimarimashita} — \textquotedblleft{}The town's population fell, so the shops closed.\textquotedblright{}
\end{quote}

\begin{cousinweb}
\textbf{\ja{じんこう}} (\emph{jinkō}) is Sino-Japanese (kango): two Chinese characters read with their on'yomi, readings borrowed from Chinese. They are \emph{jin}, \textquotedblleft{}person\textquotedblright{} (Middle Chinese \emph{nyin}) and \emph{kō}, \textquotedblleft{}mouth\textquotedblright{} (\emph{khuwX}) (KANJIDIC2; Middle Chinese in Baxter's transcription).
\end{cousinweb}

\begin{grammarlens}[title={What you've built}]
One more word for giving reasons and opinions.
\end{grammarlens}

\subsection*{Your turn}
\begin{itemize}
\raggedright
  \item \emph{Say it:} \emph{jinkō}
  \item \emph{Say it:} \emph{jinkō}, once more
  \item \emph{Say it:} \emph{machi no jinkō ga sukunaku natta node, mise ga shimarimashita} — The town's population fell, so the shops closed.
  \item \emph{Recall:} say the Japanese for young people, then the Japanese for an elderly person, then say \emph{jinkō} again
\end{itemize}

\begin{tcolorbox}[breakable,colback=teal!4,colframe=teal!35!black,title={Before you move on}]
What does \ja{じんこう} mean? (\textquotedblleft{}Population\textquotedblright{}.) Say it once more.
\end{tcolorbox}
\section[\texorpdfstring{\ja{かいごする} (kaigo suru)}{かいごする (kaigo suru)}]{\ja{かいごする} (kaigo suru) — to give care; to look after the old or sick}

\label{lesson:JA-C333-kaigo-suru}



\begin{quote}
Before the new one: say the Japanese for an elderly person, then the Japanese for population.
\end{quote}

\subsection*{You'll want to know: \ja{かいごする}}
\textbf{\ja{かいごする}} — \emph{kaigo suru} — \textquotedblleft{}to give care; to look after the old or sick\textquotedblright{}.

\begin{quote}
\textbf{\ja{そぼが} \ja{びょうきに} \ja{なったので、ははが} \ja{かいごしています。}} — \emph{sobo ga byōki ni natta node, haha ga kaigo shite imasu} — \textquotedblleft{}My grandmother fell ill, so my mother is looking after her.\textquotedblright{}
\end{quote}

Polite \textbf{\ja{かいごします}}, \emph{kaigoshimasu}; a \ja{する} verb. Past (polite): \textbf{\ja{ははが} \ja{そぼを} \ja{かいごしました。}}, \emph{haha ga sobo o kaigo shimashita}.

\begin{cousinweb}
\textbf{\ja{かいごする}} (\emph{kaigo suru}) is built from a Sino-Japanese noun and the native verb \emph{suru}, \textquotedblleft{}to do\textquotedblright{}. The noun, \emph{kaigo}, is two Chinese characters read with their on'yomi, readings borrowed from Chinese: \emph{kai}, \textquotedblleft{}mediate\textquotedblright{} (Middle Chinese \emph{keajH}) and \emph{go}, \textquotedblleft{}protect\textquotedblright{} (\emph{huH}) (KANJIDIC2; JMdict; Middle Chinese in Baxter's transcription).
\end{cousinweb}

\begin{grammarlens}[title={What you've built}]
One more word for giving reasons and opinions.
\end{grammarlens}

\subsection*{Your turn}
\begin{itemize}
\raggedright
  \item \emph{Say it:} \emph{kaigo suru}
  \item \emph{Say it:} \emph{kaigo suru}, once more
  \item \emph{Say it:} \emph{sobo ga byōki ni natta node, haha ga kaigo shite imasu} — My grandmother fell ill, so my mother is looking after her.
  \item \emph{Recall:} say the Japanese for an elderly person, then the Japanese for population, then say \emph{kaigo suru} again
\end{itemize}

\begin{tcolorbox}[breakable,colback=teal!4,colframe=teal!35!black,title={Before you move on}]
What does \ja{かいごする} mean? (\textquotedblleft{}To give care; to look after the old or sick\textquotedblright{}.) Say it once more.
\end{tcolorbox}
\section[\texorpdfstring{\ja{こそだて} (kosodate)}{こそだて (kosodate)}]{\ja{こそだて} (kosodate) — bringing up children}

\label{lesson:JA-C333-kosodate}



\begin{quote}
Before the new one: say the Japanese for population, then the Japanese for to give care.
\end{quote}

\subsection*{You'll want to know: \ja{こそだて}}
\textbf{\ja{こそだて}} — \emph{kosodate} — \textquotedblleft{}bringing up children\textquotedblright{}.

\begin{quote}
\textbf{\ja{こそだては} \ja{じかんが} \ja{かかると} \ja{おもいます。}} — \emph{kosodate wa jikan ga kakaru to omoimasu} — \textquotedblleft{}I think bringing up children takes time.\textquotedblright{}
\end{quote}

\begin{cousinweb}
\textbf{\ja{こそだて}} (\emph{kosodate}) is a native Japanese word (wago), inherited rather than borrowed from Chinese: it is built from \emph{ko}, \textquotedblleft{}child\textquotedblright{} and \emph{sodate}, \textquotedblleft{}raising (sodateru)\textquotedblright{}, and its kanji were attached later as kun'yomi (JMdict).
\end{cousinweb}

\begin{grammarlens}[title={What you've built}]
One more word for giving reasons and opinions.
\end{grammarlens}

\subsection*{Your turn}
\begin{itemize}
\raggedright
  \item \emph{Say it:} \emph{kosodate}
  \item \emph{Say it:} \emph{kosodate}, once more
  \item \emph{Say it:} \emph{kosodate wa jikan ga kakaru to omoimasu} — I think bringing up children takes time.
  \item \emph{Recall:} say the Japanese for population, then the Japanese for to give care, then say \emph{kosodate} again
\end{itemize}

\begin{tcolorbox}[breakable,colback=teal!4,colframe=teal!35!black,title={Before you move on}]
What does \ja{こそだて} mean? (\textquotedblleft{}Bringing up children\textquotedblright{}.) Say it once more.
\end{tcolorbox}
